% Minimal preamble for Open Logic Project sources (pinned commit 1e960beff9ed7835bf3e3f1335e21af3439cd107).
%
% OLP's own master preamble cannot run on the project host: TeX Live 2025 there lacks
% prooftree.sty, ntheorem.sty, extarrows.sty, mathdots.sty, isomath.sty and subfiles.sty.
% Following the project's own rule ("if the author preamble needs packages the toolchain lacks,
% do not invent macros -- use a minimal working preamble"), this file declares, verbatim in
% spirit and numbering behaviour, exactly the theorem-like environments OLP defines in
% sty/open-logic-book-envs.sty (thm/lem/prop/cor/defn/ex/prob/rem/note/case/axiom/conv),
% plus the shorthand accents OLP uses (`!!a`, `!!b`, `!!{...}`).
%
% Author text is untouched; only the packaging is ours.
\documentclass[10pt,a4paper]{article}

\usepackage{amsmath}
\usepackage{amssymb}
\usepackage{amsthm}
\usepackage{graphicx}
\usepackage{hyperref}
\usepackage{enumitem}
\usepackage[framemethod=default]{mdframed}

% ---- OLP shorthand accents (from sty/open-logic.sty) ----
% NB: \! is taken by LaTeX (negative thin space), so we only define the !!a / !!b families
% and the \olup helper used by them.
\newcommand*{\olup}[1]{\MakeUppercase{#1}}
\def\!!!a#1{\olup{#1}}
\def\!a#1{\olup{#1}}
\def\!b#1{\olup{#1}}

% ---- theorem-like environments, same family numbering as OLP ----
\theoremstyle{plain}
\newtheorem{thm}{Theorem}[section]
\newtheorem{lem}[thm]{Lemma}
\newtheorem{prop}[thm]{Proposition}
\newtheorem{cor}[thm]{Corollary}
\newtheorem{axiom}{Axiom}

\theoremstyle{definition}
\newtheorem{defn}[thm]{Definition}
\newtheorem{ex}[thm]{Example}
\newtheorem{prob}{Problem}[section]
\newtheorem{probd}{Problem}
\newtheorem{conv}[thm]{Convention}

\theoremstyle{remark}
\newtheorem{rem}[thm]{Remark}
\newtheorem{note}[thm]{Note}
\newtheorem{case}[thm]{Case}

% ---- environments OLP defines as wrappers ----
\newenvironment{probtag}[1]{\begin{prob}}{\end{prob}}
\newenvironment{explain}{}{}
\newenvironment{explain*}{}{}
\newenvironment{editorial}{}{}
\newenvironment{digress}{}{}
\newenvironment{digress*}{}{}
\newenvironment{history}{}{}
\newenvironment{reference}{}{}
\newcommand*{\olfileid}[3]{}
\newcommand*{\olsection}[1]{\section*{#1}}
\newcommand*{\olchapter}[3]{}
\newcommand*{\olpart}[2]{}
\newcommand*{\olimport}[1]{}
\newcommand*{\olinclude}[1]{}
\newcommand*{\olpath}[1]{}
\newcommand*{\ollangid}[1]{}
\newcommand*{\problemtag}[1]{}
\newcommand*{\problemauthor}[1]{}
\newcommand*{\ollabel}[1]{\label{#1}}
% \olref takes 0-3 OPTIONAL arguments plus one mandatory label:
%     \olref{lab}  |  \olref[from]{lab}  |  \olref[from][to]{lab}  |  \olref[a][b][c]{lab}
% A single-argument definition silently swallows "[sth]" and emits a broken reference
% (measured on the host: the text came out as "??iso]x"). The chain below consumes the
% brackets and cross-references only the mandatory argument.
\makeatletter
\newcommand{\olref}{\@ifnextchar[{\olref@i}{\olref@final}}
\def\olref@i[#1]{\@ifnextchar[{\olref@ii}{\olref@final}}
\def\olref@ii[#1]{\@ifnextchar[{\olref@iii}{\olref@final}}
\def\olref@iii[#1]{\olref@final}
\def\olref@final#1{\ref{#1}}
\makeatother
\newcommand*{\olnewpage}{\clearpage}

% ===================== OLP compatibility layer (1) =====================

\makeatletter

% ---------------- proof trees (OLP's prooftree is a bussproofs derivative) ----------------
\providecommand{\AxiomC}[1]{}
\providecommand{\UnaryInfC}[1]{}
\providecommand{\BinaryInfC}[1]{}
\providecommand{\TrinaryInfC}[1]{}
\providecommand{\QuaternaryInfC}[1]{}
\providecommand{\QuinaryInfC}[1]{}
\providecommand{\RightLabel}[1]{}
\providecommand{\LeftLabel}[1]{}
\providecommand{\doubleLine}{}
\providecommand{\noLine}{}
\providecommand{\DisplayProof}{}
\providecommand{\fCenter}{}
\newenvironment{prooftree}{}{}
\newenvironment{ltprooftree}{}{}
\newenvironment{oltableau}{}{}

% ---------------- logical notation ----------------
\providecommand{\lif}{\mathbin{\rightarrow}}
\providecommand{\lforall}{\mathord{\forall}}
\providecommand{\lexists}{\mathord{\exists}}
\providecommand{\ltrue}{\top}
\providecommand{\lfalse}{\bot}
\providecommand{\Ltrue}{\top}
\providecommand{\Lfalse}{\bot}
\providecommand{\Sequent}{\mathrel{\Rightarrow}}
\providecommand{\Proves}{\mathrel{\vdash}}
\providecommand{\Entails}{\mathrel{\models}}
\providecommand{\Subst}[3]{#1[#2/#3]}
\providecommand{\sFmla}[2]{#1(#2)}
\providecommand{\pFmla}[2]{#1(#2)}
\providecommand{\formula}{\relax}
\providecommand{\True}{\top}
\providecommand{\False}{\bot}
\providecommand{\Args}[1]{#1}
\providecommand{\Ts}{\ensuremath{\top}}
\providecommand{\Fi}{\ensuremath{\bot}}
\providecommand{\lfalsum}{\bot}

% ---------------- objects, structures, metatheory ----------------
\providecommand{\Obj}{}
\providecommand{\Struct}{\ensuremath{\mathfrak{A}}}
\providecommand{\Domain}{\ensuremath{|\mathfrak{A}|}}
\providecommand{\Assign}{\ensuremath{s}}
\providecommand{\Sat}[2]{\ensuremath{#1 \models #2}}
\providecommand{\mSat}[3]{\ensuremath{#1 \models #2[#3]}}
\providecommand{\mModel}[2]{\ensuremath{#1 \models #2}}
\providecommand{\Lang}{\ensuremath{\mathcal{L}}}
\providecommand{\Th}{\ensuremath{\mathrm{Th}}}
\providecommand{\Log}{\ensuremath{\mathrm{L}}}
\providecommand{\tuple}[1]{\langle #1 \rangle}
\providecommand{\num}[1]{\underline{#1}}
\providecommand{\cfind}[1]{\varphi_{#1}}
\providecommand{\Nat}{\ensuremath{\mathbb{N}}}
\providecommand{\Mult}{\ensuremath{\mathrm{Mult}}}
\providecommand{\Add}{\ensuremath{\mathrm{Add}}}
\providecommand{\Exp}{\ensuremath{\mathrm{Exp}}}
\providecommand{\Zero}{\ensuremath{0}}
\providecommand{\Succ}{\ensuremath{\mathrm{S}}}
\providecommand{\lambd}{\ensuremath{\lambda}}
\providecommand{\bredpar}{\mathrel{\rightarrow_\beta}}
\providecommand{\bcd}[1]{\ensuremath{\ulcorner #1 \urcorner}}
\providecommand{\Set}[1]{\ensuremath{\{#1\}}}
\providecommand{\indcase}[1]{\par\noindent\emph{Case #1}.\quad}
\providecommand{\indfrm}[1]{\par\noindent\emph{Form #1}.\quad}
\providecommand{\iftag}[1]{#1}
\providecommand{\tagitem}[1]{\item[#1]}
\providecommand{\Elim}[1]{}
\providecommand{\Intro}[1]{}
\providecommand{\LeftR}[1]{}
\providecommand{\RightR}[1]{}
\providecommand{\TAss}[1]{}
\providecommand{\TRule}[1]{}

\makeatother
\providecommand{\olpath}{.}

% ===================== OLP compatibility layer (2) =====================
% --- second compatibility block: notation and shorthands used by the frozen OLP units ---
\providecommand{\CutCS}{\mathrm{Cut}}
\providecommand{\LeftSubproofLabel}{\relax}
\providecommand{\RightSubproofLabel}{\relax}
\providecommand{\maeh}{\relax}
\providecommand{\cheight}{\relax}
\providecommand{\ande}{\mathrm{and}_E}
\providecommand{\andi}{\mathrm{and}_I}
\providecommand{\ore}{\mathrm{or}_E}
\providecommand{\ori}{\mathrm{or}_I}
\providecommand{\typeof}{\mathord{:}}
\providecommand{\pValue}[1]{#1}
\providecommand{\Sent}{\mathrm{Sent}}
\providecommand{\cardplus}{\mathbin{+}}
\providecommand{\cardtimes}{\mathbin{\cdot}}
\providecommand{\cardexpo}{\mathbin{\hat{\ }}}
\providecommand{\cardnless}{\not<}
\providecommand{\cardle}{\leq}
\providecommand{\funfromto}[3]{#1\colon #2 \to #3}
\providecommand{\funimage}[1]{\operatorname{im}(#1)}
\providecommand{\fdefined}{\mathord{\downarrow}}
\providecommand{\canonord}{\operatorname{ord}}
\providecommand{\ordtimes}{\mathbin{\cdot}}
\providecommand{\ordeq}[2]{#1 = #2}
\providecommand{\ordneq}[2]{#1 \neq #2}
\providecommand{\ordtype}[1]{\operatorname{ord}(#1)}
\providecommand{\ordsucc}{\mathord{+1}}
\providecommand{\ordunion}{\cup}
\providecommand{\onesphere}{S^1}
\providecommand{\rotationsgroup}{\mathrm{SO}(2)}
\providecommand{\equivrep}[1]{[#1]}
\providecommand{\rlexless}{\prec}
\providecommand{\ZFC}{\mathrm{ZFC}}
\providecommand{\ZF}{\mathrm{ZF}}
\providecommand{\Rat}{\mathbb{Q}}
\providecommand{\Real}{\mathbb{R}}
\providecommand{\PosInt}{\mathbb{Z}^{+}}
\providecommand{\Id}{\mathrm{id}}
\providecommand{\Bin}{\mathrm{Bin}}
\providecommand{\pto}{\rightharpoonup}
\providecommand{\dom}{\operatorname{dom}}
\providecommand{\ran}{\operatorname{ran}}
\providecommand{\len}{\operatorname{len}}
\providecommand{\concat}{\mathbin{{}^\frown}}
\providecommand{\card}[1]{\lvert #1 \rvert}
\providecommand{\Pow}{\mathcal{P}}
\providecommand{\Setabs}[2]{\{\, #1 \mid #2 \,\}}
\providecommand{\comp}{\circ}
\providecommand{\Char}{\chi}
\providecommand{\Prf}{\operatorname{Prf}}
\providecommand{\Prov}{\operatorname{Prov}}
\providecommand{\Proj}{\operatorname{proj}}
\providecommand{\Gn}{\operatorname{gn}}
\providecommand{\gn}{\operatorname{gn}}
\providecommand{\Value}[1]{\operatorname{val}(#1)}
\providecommand{\Deduce}{\mathrm{D}}
\providecommand{\Dual}{\mathrm{d}}
\providecommand{\Ax}{\mathrm{Ax}}
\providecommand{\pSat}[3]{#1 \models #2\,[#3]}
\providecommand{\pAssign}{\sigma}
\providecommand{\bforall}{\forall}
\providecommand{\bexists}{\exists}
\providecommand{\liff}{\mathbin{\leftrightarrow}}
\providecommand{\FV}{\operatorname{FV}}
\providecommand{\UnaryInf}{\ensuremath{\mid}}
\providecommand{\BinaryInf}{\ensuremath{\mid}}
\providecommand{\depth}{\operatorname{depth}}
\providecommand{\oliflabeldef}[3]{#1}
\providecommand{\TMblank}{\sqcup}
\providecommand{\TMstroke}{\mid}
\providecommand{\TMstay}{\mathrm{S}}
\providecommand{\TMright}{\mathrm{R}}
\providecommand{\TMleft}{\mathrm{L}}
\providecommand{\TMendtape}{\triangleright}
\providecommand{\umin}{\mu}
\providecommand{\Sum}{\sum}
\providecommand{\citep}[1]{\cite{#1}}
\providecommand{\citet}[1]{\cite{#1}}
\providecommand{\citeauthor}[1]{\cite{#1}}
\providecommand{\citeyear}[1]{\cite{#1}}

% ===================== OLP compatibility layer (3) =====================
% --- third compatibility block: remaining notation seen in the frozen OLP units ---
\providecommand{\iso}{\cong}
\providecommand{\nszero}{\mathrm{0}}
\providecommand{\nssucc}{\mathrm{S}}
\providecommand{\nsplus}{\mathbin{+}}
\providecommand{\nssless}{<}
\providecommand{\nsless}{<}
\providecommand{\PL}{\mathrm{PL}}
\providecommand{\Nec}{\mathord{\Box}}
\providecommand{\RK}{\mathrm{RK}}
\providecommand{\ST}{\mathrm{ST}}
\providecommand{\SSubst}[3]{#1[#2/#3]}
\providecommand{\subst}[3]{#1[#2/#3]}
\providecommand{\tagfalse}{\mathrm{F}}
\providecommand{\cutr}{\mathrm{Cut}}
\providecommand{\cardneq}{\neq}
\providecommand{\issucc}[1]{\operatorname{succ}(#1)}
\providecommand{\tup}[1]{\langle #1 \rangle}
\providecommand{\Var}{\mathrm{Var}}
\providecommand{\Bew}{\operatorname{Bew}}
\providecommand{\Con}{\operatorname{Con}}
\providecommand{\Cons}{\operatorname{Cons}}
\providecommand{\pDeduce}{\mathrm{D}}
\providecommand{\psat}{\models}
\providecommand{\den}[1]{\llbracket #1 \rrbracket}
\providecommand{\Sem}[2]{\llbracket #1 \rrbracket_{#2}}
\providecommand{\DomainOf}[1]{\lvert #1 \rvert}
\providecommand{\Form}{\mathrm{Form}}
\providecommand{\Terms}{\mathrm{Term}}
\providecommand{\Atom}{\mathrm{Atom}}
\providecommand{\FVars}[1]{\operatorname{FV}(#1)}
\providecommand{\fv}[1]{\operatorname{fv}(#1)}
\providecommand{\tuplep}[1]{\langle #1 \rangle}
\providecommand{\restrict}[2]{#1 \upharpoonright #2}
\providecommand{\sto}{\rightarrow}
\providecommand{\FmlaSet}[1]{\mathrm{Fmla}_{#1}}
\providecommand{\Thm}{\mathrm{Thm}}
\providecommand{\Eval}{\operatorname{val}}
\providecommand{\blank}{\sqcup}
\providecommand{\blankTape}{\sqcup}
\providecommand{\Consis}{\operatorname{Con}}
\providecommand{\drs}{\mathrm{d}}
\providecommand{\VRule}{\mid}
\providecommand{\Bet}{\mathrm{B}}
\providecommand{\Seg}[1]{\mathrm{Seg}_{#1}}
\providecommand{\Enums}{\mathrm{En}}
\providecommand{\Z}{\mathbb{Z}}
\providecommand{\N}{\mathbb{N}}
\providecommand{\abs}[1]{\lvert #1 \rvert}
\providecommand{\floor}[1]{\lfloor #1 \rfloor}
\providecommand{\ceil}[1]{\lceil #1 \rceil}

% ===================== OLP compatibility layer (4) =====================
% --- fourth compatibility block: harvested from the per-candidate compile logs ---
\providecommand{\beredpar}{\mathrel{\rightarrow_\beta}}
\makeatletter
\def\bered{\@ifnextchar[{\bered@i}{\mathrel{\rightarrow_\beta}}}
\def\bered@i[#1]{\mathrel{\rightarrow_\beta}}
\def\bred{\@ifnextchar[{\bred@i}{\mathrel{\rightarrow_\beta}}}
\def\bred@i[#1]{\mathrel{\rightarrow_\beta}}
\makeatother
\providecommand{\beredone}{\mathrel{=_\beta}}
\providecommand{\bredone}{\mathrel{=_\beta}}
\providecommand{\xredone}{\mathrel{=_\beta}}
\providecommand{\aconvone}{\mathrel{=_\alpha}}
\providecommand{\cardless}{<}
\providecommand{\tagrefs}[1]{#1}
\providecommand{\ordplus}{\mathbin{+}}
\providecommand{\funrestrictionto}[2]{#1 \upharpoonright #2}
\providecommand{\fundefined}{\mathord{\uparrow}}
\providecommand{\Trm}{\mathrm{Trm}}
\providecommand{\OProv}{\operatorname{Prov}}
\providecommand{\OPrf}{\operatorname{Prf}}
\providecommand{\aeq}{\mathrel{=}}
\providecommand{\mClass}[1]{[#1]}
\providecommand{\equivclass}[1]{[#1]}
\providecommand{\pheight}{\operatorname{ht}}
\providecommand{\phy}{\operatorname{ht}}
\providecommand{\Cut}{\mathrm{Cut}}
\providecommand{\QuantRank}{\operatorname{qr}}
\providecommand{\elemequiv}{\mathrel{\equiv}}
\providecommand{\PIso}{\mathrm{PIso}}
\providecommand{\Theory}{\mathrm{Th}}
\providecommand{\supstrict}{\subsetneq}
\providecommand{\trcl}{\operatorname{trcl}}
\providecommand{\cref}[1]{\ref{#1}}
\providecommand{\Discharge}{\mathrm{Dis}}
\providecommand{\Weakening}{\mathrm{Wk}}
\providecommand{\ZFminus}{\mathrm{ZF}^{-}}
\makeatletter
\renewcommand{\Dual}{\@ifnextchar[{\Dual@i}{\ifmmode\mathrm{d}\else$\mathrm{d}$\fi}}
\def\Dual@i[#1]{\ifmmode\mathrm{d}[#1]\else$\mathrm{d}[#1]$\fi}
\makeatother
\providecommand{\RightR}{\mathrm{R}}
\providecommand{\LeftR}{\mathrm{L}}
\providecommand{\Intro}{\mathrm{I}}
\providecommand{\Axiom}{\mathrm{Ax}}
\providecommand{\formula}[1]{#1}
\providecommand{\stared}{\ast}
\providecommand{\oset}[2]{\{\, #1 \mid #2 \,\}}
\providecommand{\setof}[2]{\{\, #1 \mid #2 \,\}}
\providecommand{\pow}[1]{\mathcal{P}(#1)}
\providecommand{\ordpair}[2]{\langle #1, #2 \rangle}
\providecommand{\ordtriple}[3]{\langle #1, #2, #3 \rangle}
\providecommand{\union}{\cup}
\providecommand{\inter}{\cap}
\providecommand{\seq}[1]{\langle #1 \rangle}
\providecommand{\SuchThat}{\mid}
\providecommand{\where}{\mid}
\providecommand{\TVar}{\mathrm{TVar}}
\providecommand{\Bool}{\mathrm{Bool}}
\providecommand{\Val}{\operatorname{Val}}
\providecommand{\TL}[1]{\mathrm{TL}}
\providecommand{\xstared}{\ast}

% ===================== OLP compatibility layer (5) =====================
% --- fifth compatibility block: rescues the quarantined OLP candidates ---
\providecommand{\ident}{\operatorname{id}}
\providecommand{\setrank}{\operatorname{rank}}
\providecommand{\cardfont}{\mathfrak{c}}
\providecommand{\Complement}{\overline}
\providecommand{\QR}{\mathrm{QR}}
\providecommand{\becd}{\ensuremath{\ulcorner}}
\providecommand{\disjointsum}{\mathbin{\oplus}}
\providecommand{\cardeq}[2]{#1 = #2}
\providecommand{\xred}{\mathrel{\rightarrow}}
% --- sixth compatibility block: names reported by real compile logs ---
\providecommand{\ORefut}{\operatorname{Refut}}
\providecommand{\ORProv}{\operatorname{Prov}}
\providecommand{\cutrank}{\operatorname{cr}}
\providecommand{\Refut}{\operatorname{Refut}}
\providecommand{\DischargeRule}{\mathrm{Dis}}
\providecommand{\Contraction}{\mathrm{C}}
\providecommand{\bottomAlignProof}{\relax}
\providecommand{\FalseCl}{\bot}
\providecommand{\red}{\mathrel{\rightarrow}}
\providecommand{\varAssign}{\sigma}
\providecommand{\defiff}{\mathrel{\leftrightarrow}}
\providecommand{\FalseInt}{\bot}
\providecommand{\fn}{\operatorname{fn}}
\providecommand{\dotted}{\cdot}

% --- tag-macro semantic fix (2026-10-13, final): \iftag CHOOSES a branch and the OLP proof body is
% always the SECOND group (\iftag{tag}{body}, {tag}{body}{c}, {tag}{body}{c}{d}); a definition that
% prints the tag literal drops the proof body. Bodies contain blank lines, hence \long.
\makeatletter
\providecommand{\iftag}{}
\long\def\iftag#1#2{\@ifnextchar\bgroup{\iftag@three{#2}}{\iftag@stop{#2}}}
\long\def\iftag@three#1#2{\iftag@stop{#1}\@ifnextchar\bgroup{\iftag@eat}{\relax}}
\long\def\iftag@eat#1{}
\long\def\iftag@stop#1{#1}
\long\def\oliflabeldef#1#2{\@ifnextchar\bgroup{\@gobble}{\iftag@stop{#2}}}
\providecommand{\tagrefs}[1]{}
\providecommand{\settag}[1]{}
\makeatother
\begin{document}

% ==== OLP file: content/computability/recursive-functions/bounded-minimization.tex ====
% Part: computability
% Chapter: recursive-functions
% Section: bounded-minimization



\olfileid{cmp}{rec}{bmi}
\olsection{Bounded Minimization}

\begin{explain}
It is often useful to define a function as the least number satisfying
some property or relation~$P$. If $P$ is decidable, we can compute
this function simply by trying out all the possible numbers, $0$, $1$,
$2$, \dots, until we find the least one satisfying~$P$.  This kind of
unbounded search takes us out of the realm of primitive recursive
functions. However, if we're only interested in the least number
\emph{less than some independently given bound}, we stay primitive
recursive. In other words, and a bit more generally, suppose we have a
primitive recursive relation~$R(x,z)$. Consider the function that maps
$x$ and~$y$ to the least $z < y$ such that $R(x, z)$. It, too, can be
computed, by testing whether $R(x, 0)$, $R(x, 1)$, \dots, $R(x, y-1)$.
But why is it primitive recursive?
\end{explain}

\begin{prop}
If $R(\vec x, z)$ is primitive recursive, so is the function
$m_R(\vec{x}, y)$ which returns the least~$z$ less than~$y$ such that
$R(\vec x, z)$ holds, if there is one, and $y$ otherwise.  We will
write the function~$m_R$ as
\[
\bmin{z < y}{R(\vec{x}, z)},
\]
\end{prop}

\begin{proof}
Note than there can be no~$z < 0$ such that $R(\vec{x}, z)$ since
there is no $z < 0$ at all.  So $m_R(\vec x, 0) = 0$.

In case the bound is of the form $y + 1$ we have three cases:
\begin{enumerate}
\item There is a $z < y$ such that $R(\vec{x}, z)$, in which case
$m_R(\vec{x}, y+1) = m_R(\vec{x}, y)$.
\item There is no such~$z<y$ but $R(\vec{x}, y)$ holds, then
$m_R(\vec{x}, y+1) = y$.
\item There is no $z < y+1$ such that $R(\vec{x}, z)$, then
$m_R(\vec{z}, y+1) = y+1$.
\end{enumerate}
So we can define $m_R(\vec x, 0)$ by primitive recursion as follows:
\begin{align*}
m_R(\vec{x}, 0) & = 0\\
m_R(\vec{x}, y+1) & =
\begin{cases}
m_R(\vec{x}, y) & \text{if $m_R(\vec x, y) \neq y$}\\
y & \text{if $m_R(\vec x, y) = y$ and $R(\vec{x}, y)$}\\
y+1 & \text{otherwise.}
\end{cases}
\end{align*}
Note that there is a $z<y$ such that $R(\vec x, z)$ iff $m_R(\vec x,
y) \neq y$.
\end{proof}

\begin{prob}
Suppose $R(\vec x, z)$ is primitive recursive. Define the function
$m'_R(\vec{x}, y)$ which returns the least~$z$ less than~$y$ such that
$R(\vec{x}, z)$ holds, if there is one, and $0$ otherwise, by
primitive recursion from~$\Char{R}$.
\end{prob}


% ==== OLP file: content/computability/recursive-functions/composition.tex ====
% Part: computability
% Chapter: recursive-functions
% Section: composition



\olfileid{cmp}{rec}{com}
\olsection{Composition}

If $f$ and $g$ are two one-place functions of natural numbers, we can
compose them: $h(x) = f(g(x))$. The new function~$h(x)$ is then
defined by \emph{composition} from the functions $f$ and~$g$. We'd
like to generalize this to functions of more than one argument.

Here's one way of doing this: suppose $f$ is a $k$-place function,
and $g_0$, \dots, $g_{k-1}$ are $k$ functions which are all
$n$-place. Then we can define a new $n$-place function $h$ as follows:
\[
h(x_0, \dots, x_{n-1}) =
f(g_0(x_0, \dots, x_{n-1}), \dots, g_{k-1}(x_0, \dots, x_{n-1}))
\]
If $f$ and all $g_i$ are computable, so is $h$: To compute $h(x_0,
\dots, x_{n-1})$, first compute the values $y_i = g_i(x_0, \dots,
x_{n-1})$ for each $i = 0$, \dots,~$k-1$. Then feed these values into
$f$ to compute $h(x_0, \dots, x_{k-1}) = f(y_0, \dots, y_{k-1})$.

This may seem like an overly restrictive characterization of what
happens when we compute a new function using some existing ones. For
one thing, sometimes we do not use all the arguments of a function, as
when we defined $g(x, y, z) = \Succ(z)$ for use in the primitive
recursive definition of~$\Add$. Suppose we are allowed use of the
following functions:
\[
\Proj{n}{i}(x_0, \dots, x_{n-1}) = x_i
\]
The functions~$\Proj{k}{i}$ are called \emph{projection} functions:
$\Proj{n}{i}$ is an $n$-place function. Then $g$ can be defined by
\[
g(x, y, z) = \Succ(\Proj{3}{2}(x, y, z)).
\]
Here the role of $f$ is played by the $1$-place function $\Succ$, so
$k=1$. And we have one $3$-place function $\Proj{3}{2}$ which plays
the role of $g_0$. The result is a $3$-place function that returns the
successor of the third argument.

The projection functions also allow us to define new functions by
reordering or identifying arguments. For instance, the function $h(x)
= \Add(x, x)$ can be defined by
\[
h(x_0) = \Add(\Proj{1}{0}(x_0),\Proj{1}{0}(x_0)).
\]
Here $k=2$, $n=1$, the role of $f(y_0,y_1)$ is played by $\Add$, and
the roles of $g_0(x_0)$ and $g_1(x_0)$ are both played
by~$\Proj{1}{0}(x_0)$, the one-place projection function (aka the
identity function).

If $f(y_0, y_1)$ is a function we already have, we can define the
function $h(x_0, x_1) = f(x_1, x_0)$ by
\[
h(x_0, x_1) = f(\Proj{2}{1}(x_0, x_1),\Proj{2}{0}(x_0, x_1)).
\]
Here $k=2$, $n = 2$, and the roles of $g_0$ and $g_1$ are played by
$\Proj{2}{1}$ and~$\Proj{2}{0}$, respectively.

You may also worry that $g_0$, \dots,~$g_{k-1}$ are all required
to have the same arity~$n$. (Remember that the \emph{arity} of a
function is the number of arguments; an $n$-place function has
arity~$n$.) But adding the projection functions provides the desired
flexibility. For example, suppose $f$ and~$g$ are $3$-place functions
and $h$~is the $2$-place function defined by
\[
h(x,y) = f(x,g(x,x,y),y).
\]
The definition of~$h$ can be rewritten with the projection
functions, as
\[
h(x,y) = f(\Proj{2}{0}(x,y), g(\Proj{2}{0}(x,y), \Proj{2}{0}(x,y),
\Proj{2}{1}(x,y)), \Proj{2}{1}(x,y)).
\]
Then $h$ is the composition of $f$ with $\Proj{2}{0}$, $l$, and
$\Proj{2}{1}$, where
\[
l(x,y) = g(\Proj{2}{0}(x,y),\Proj{2}{0}(x,y),\Proj{2}{1}(x,y)),
\]
i.e., $l$ is the composition of $g$ with $\Proj{2}{0}$, $\Proj{2}{0}$,
and~$\Proj{2}{1}$.


% ==== OLP file: content/computability/recursive-functions/examples.tex ====
% Part: computability
% Chapter: recursive-functions
% Section: examples



\olfileid{cmp}{rec}{exa}
\olsection{Examples of Primitive Recursive Functions}


We already have some examples of primitive recursive functions: the
addition and multiplication functions~$\Add$ and $\Mult$.  The
identity function $\fn{id}(x) = x$ is primitive recursive, since it is
just~$\Proj{1}{0}$. The constant functions $\fn{const}_n(x) = n$ are
primitive recursive since they can be defined from $\Zero$ and $\Succ$
by successive composition. This is useful when we want to use
constants in primitive recursive definitions, e.g., if we want to
define the function $f(x) = 2 \cdot x$ can obtain it by composition
from $\fn{const}_n(x)$ and multiplication as $f(x) =
\Mult(\fn{const}_2(x), \Proj{1}{0}(x))$. We'll make use of this trick
from now on.

\begin{prop}
  The exponentiation function $\fn{exp}(x, y) = x^y$ is primitive recursive.
\end{prop}

\begin{proof}
  We can define $\fn{exp}$ primitive recursively as
  \begin{align*}
    \fn{exp}(x, 0) & = 1\\
    \fn{exp}(x, y+1) & = \Mult(x, \fn{exp}(x,y)).
    \intertext{Strictly speaking, this is not a recursive definition
      from primitive recursive functions. Officially, though, we
      have:}
    \fn{exp}(x, 0) & = f(x)\\
    \fn{exp}(x, y+1) & = g(x, y, \fn{exp}(x,y)).
    \intertext{where}
    f(x) & = \Succ(\Zero(x)) = 1\\
    g(x, y, z) & = \Mult(\Proj{3}{0}(x, y, z), \Proj{3}{2}(x, y, z)) = x \cdot z
  \end{align*}
  and so $f$ and $g$ are defined from primitive recursive functions by
  composition.
\end{proof}

\begin{prop}
  The predecessor function $\fn{pred}(y)$ defined by
  \[
  \fn{pred}(y) = \begin{cases}
    0 & \text{if $y=0$}\\
    y-1 & \text{otherwise}
  \end{cases}
  \]
  is primitive recursive.
\end{prop}

\begin{proof}
 Note that 
 \begin{align*}
   \fn{pred}(0) & = 0 \text{ and}\\
   \fn{pred}(y+1) & = y.
 \end{align*}
 This is almost a primitive recursive definition.  It does not,
 strictly speaking, fit into the pattern of definition by primitive
 recursion, since that pattern requires at least one extra
 argument~$x$. It is also odd in that it does not actually use
 $\fn{pred}(y)$ in the definition of $\fn{pred}(y+1)$. But we can
 first define $\fn{pred}'(x, y)$ by
 \begin{align*}
   \fn{pred}'(x, 0) & = \Zero(x) = 0,\\
   \fn{pred}'(x, y+1) & = \Proj{3}{1}(x, y, \fn{pred'}(x, y)) = y.
 \end{align*}
and then define $\fn{pred}$ from it by composition, e.g., as
$\fn{pred}(x) = \fn{pred}'(\Zero(x), \Proj{1}{0}(x))$.
\end{proof}

\begin{prop}
  The factorial function $\fn{fac}(x) = \fact{x} = 1 \cdot 2 \cdot 3
  \cdot \dots \cdot x$ is primitive recursive.
\end{prop}

\begin{proof}
  The obvious primitive recursive definition is
  \begin{align*}
    \fn{fac}(0) &= 1\\
    \fn{fac}(y+1) & = \fn{fac}(y) \cdot (y+1).
    \intertext{Officially, we have to first define a two-place function $h$}
    h(x, 0) & = \fn{const}_1(x)\\
    h(x, y+1) & = g(x, y, h(x, y))
    \intertext{where $g(x, y, z) = \Mult(\Proj{3}{2}(x, y, z),
      \Succ(\Proj{3}{1}(x, y, z)))$ and then let}
    \fn{fac}(y) & = h(\Proj{1}{0}(y), \Proj{1}{0}(y)) = h(y,y).
  \end{align*}
  From now on we'll be a bit more laissez-faire and not give the official
  definitions by composition and primitive recursion.
\end{proof}

\begin{prop}
  Truncated subtraction, $x \tsub y$, defined by
  \[
  x \tsub y = \begin{cases}
    0 & \text{if $x < y$}\\
    x-y & \text{otherwise}
  \end{cases}
  \]
  is primitive recursive.
\end{prop}

\begin{proof}
  We have:
  \begin{align*}
    x \tsub 0 & = x\\
    x \tsub (y+1) & = \fn{pred}(x \tsub y) 
  \end{align*}
\end{proof}

\begin{prop}
 The distance between $x$ and $y$, $\left|x-y\right|$, is primitive recursive.
\end{prop}

\begin{proof}
  We have $\left| x-y \right| = (x \tsub y) + (y \tsub x)$, so
  the distance can be defined by composition from $+$ and $\tsub$,
  which are primitive recursive.
\end{proof}

\begin{prop}
  The maximum of $x$ and $y$, $\fn{max}(x,y)$, is primitive recursive.
\end{prop}

\begin{proof}
  We can define $\fn{max}(x,y)$ by composition from $+$ and $\tsub$ by
  \[
  \fn{max}(x,y) \defis x + (y \tsub x).
  \]
  If $x$ is the maximum, i.e., $x \ge y$, then $y \tsub x = 0$, so $x
  + (y \tsub x) = x + 0 = x$. If $y$ is the maximum, then $y \tsub x =
  y - x$, and so $x + (y \tsub x) = x + (y - x) = y$.
\end{proof}

\begin{prop}
  \ollabel{prop:min-pr}
  The minimum of $x$ and $y$, $\fn{min}(x,y)$, is primitive recursive.
\end{prop}

\begin{proof}
  Exercise.
\end{proof}

\begin{prob}
  Prove \olref[cmp][rec][exa]{prop:min-pr}.
\end{prob}

\begin{prob}
Show that \[
f(x, y) =
2^{(2^{\iddots^{2^{x}}})}\raisebox{1ex}{\bigg\rbrace}
\raisebox{1ex}{\text {$y$ $2$'s}}\] is primitive recursive.
\end{prob}

\begin{prob}
Show that integer division $d(x, y) = \lfloor x/y \rfloor$ (i.e.,
division, where you disregard everything after the decimal point) is
primitive recursive. When $y = 0$, we stipulate $d(x, y) = 0$. Give an
explicit definition of~$d$ using primitive recursion and
composition.
\end{prob}

\begin{prop}
The set of primitive recursive functions is closed under the
following two operations:
\begin{enumerate}
\item Finite sums: if $f(\vec x, z)$ is primitive recursive, then so
is the function
\[
g(\vec x, y) \defis \sum_{z = 0}^y f(\vec x, z).
\]
\item Finite products: if $f(\vec x, z)$ is primitive recursive, then
so is the function
\[
h(\vec x, y) \defis \prod_{z = 0}^y f(\vec x, z).
\]
\end{enumerate}
\end{prop}

\begin{proof}
For example, finite sums are defined recursively by the equations
\begin{align*}
  g(\vec x, 0) & = f(\vec x, 0)\\
  g(\vec x, y+1) & = g(\vec x, y) + f(\vec x, y+1).
\end{align*}
\end{proof}



% ==== OLP file: content/computability/recursive-functions/general-recursive-functions.tex ====
% Part: computability
% Chapter: recursive-functions
% Section: general-recursive-functions



\olfileid{cmp}{rec}{gen}
\olsection{General Recursive Functions}

There is another way to obtain a set of total functions. Say a total
function $f(x,\vec z)$ is \emph{regular} if for every sequence of
natural numbers $\vec z$, there is an $x$ such that $f(x,\vec z) = 0$.
In other words, the regular functions are exactly those functions to
which one can apply unbounded search, and end up with a total
function. One can, conservatively, restrict unbounded search to
regular functions:

\begin{defn}
\ollabel{defn:general-recursive}
The set of \emph{general recursive functions} is the smallest set of
functions from the natural numbers to the natural numbers (of various
arities) containing zero, successor, and projections, and closed under
composition, primitive recursion, and unbounded search applied to
\emph{regular} functions.
\end{defn}

Clearly every general recursive function is total. The difference
between \olref{defn:general-recursive} and
\olref[par]{defn:recursive-fn} is that in the latter one is allowed to
use partial recursive functions along the way; the only requirement is
that the function you end up with at the end is total. So the word
``general,'' a historic relic, is a misnomer; on the surface,
\olref{defn:general-recursive} is \emph{less} general than
\olref[par]{defn:recursive-fn}. But, fortunately, the difference is
illusory; though the definitions are different, the set of general
recursive functions and the set of recursive functions are one and the
same.



% ==== OLP file: content/computability/recursive-functions/halting-problem.tex ====
% Part: computability
% Chapter: recursive-functions
% Section: halting-problem



\olfileid{cmp}{rec}{hlt}
\olsection{The Halting Problem}

The \emph{halting problem} in general is the problem of deciding,
given the specification~$e$ (e.g., program) of a computable function
and a number~$n$, whether the computation of the function on input~$n$
halts, i.e., produces a result.  Famously, Alan Turing proved that
this problem itself cannot be solved by a computable function, i.e.,
the function
\[
h(e, n) =
\begin{cases}
1 & \text{if computation $e$ halts on input $n$}\\
0 & \text{otherwise,}
\end{cases}
\]
is not computable.

In the context of partial recursive functions, the role of the
specification of a program may be played by the index~$e$ given in
Kleene's normal form theorem.  If $f$ is a partial recursive function,
any $e$~for which the equation in the normal form theorem holds, is an
index of~$f$. Given a number~$e$, the normal form theorem states that
\[
\cfind{e}(x) \simeq U(\mu s \; T(e, x, s))
\]
is partial recursive, and for every partial recursive $f\colon \Nat
\to \Nat$, there is an $e \in \Nat$ such that $\cfind{e}(x) \simeq
f(x)$ for all~$x \in \Nat$.  In fact, for each such $f$ there is not
just one, but infinitely many such~$e$.  The \emph{halting function}~$h$
is defined by
\[
h(e, x) =
\begin{cases}
1 & \text{if $\cfind{e}(x) \fdefined$}\\
0 & \text{otherwise.}
\end{cases}
\]
Note that $h(e, x) = 0$ if $\cfind{e}(x) \fundefined$, but also
when~$e$ is not the index of a partial recursive function at all.

\begin{thm}
\ollabel{thm:halting-problem}
The halting function~$h$ is not partial recursive.
\end{thm}

\begin{proof}
If $h$
were partial recursive, we could define
\[
d(y) =
\begin{cases}
1 & \text{if $h(y, y) = 0$}\\
\umin{x}{x \neq x} & \text{otherwise.}
\end{cases}
\]
Since no number $x$ satisfies $x \neq x$, there is no $\umin{x}{x \neq
x}$, and so $d(y) \fundefined$ iff $h(y,y) \neq 0$. From this
definition it follows that
\begin{enumerate}
\item $d(y) \fdefined$ iff $\cfind{y}(y) \fundefined$ or $y$ is
  not the index of a partial recursive function.
\item $d(y) \fundefined$ iff $\cfind{y}(y) \fdefined$.
\end{enumerate}
If $h$ were partial recursive, then $d$ would be partial recursive as
well.  Thus, by the Kleene normal form theorem, it has an index~$e_d$.
Consider the value of $h(e_d, e_d)$. There are two possible cases, $0$
and~$1$.
\begin{enumerate}
\item If $h(e_d, e_d) = 1$ then $\cfind{e_d}(e_d) \fdefined$.  But
  $\cfind{e_d} \simeq d$, and $d(e_d)$ is defined iff $h(e_d, e_d) = 0$.
  So $h(e_d, e_d) \neq 1$.
\item If $h(e_d, e_d) = 0$ then either $e_d$ is not the index of a
  partial recursive function, or it is and $\cfind{e_d}(e_d)
  \fundefined$. But again, $\cfind{e_d} \simeq d$, and $d(e_d)$ is undefined
  iff $\cfind{e_d}(e_d) \fdefined$.
\end{enumerate}
The upshot is that $e_d$ cannot, after all, be the index of a partial
recursive function.  But if $h$ were partial recursive, $d$ would be
too, and so our definition of $e_d$ as an index of it would be
admissible.  We must conclude that $h$ cannot be partial recursive.
\end{proof}




% ==== OLP file: content/computability/recursive-functions/introduction.tex ====
% Part: computability
% Chapter: recursive-functions
% Section: introduction



\olfileid{cmp}{rec}{int}
\olsection{Introduction}

In order to develop a mathematical theory of computability, one has to,
first of all, develop a \emph{model} of computability.  We now think of
computability as the kind of thing that computers do, and computers
work with symbols.  But at the beginning of the development of
theories of computability, the paradigmatic example of computation was
\emph{numerical} computation.  Mathematicians were always interested
in number-theoretic functions, i.e., functions $f\colon \Nat^n \to
\Nat$ that can be computed. So it is not surprising that at the
beginning of the theory of computability, it was such functions that
were studied.  The most familiar examples of computable numerical
functions, such as addition, multiplication, exponentiation (of
natural numbers) share an interesting feature: they can be defined
\emph{recursively}.  It is thus quite natural to attempt a general
definition of \emph{computable function} on the basis of recursive
definitions.  Among the many possible ways to define number-theoretic
functions recursively, one particularly simple pattern of definition
here becomes central: so-called \emph{primitive recursion}.

In addition to computable functions, we might be interested in
computable sets and relations. A set is computable if we can compute
the answer to whether or not a given number is \olup{element} of the
set, and a relation is computable iff we can compute whether or not a
tuple $\tuple{n_1, \dots, n_k}$ is \olup{element} of the relation.  By
considering the \emph{characteristic function} of a set or relation,
discussion of computable sets and relations can be subsumed under that
of computable functions.  Thus we can define primitive recursive
relations as well, e.g., the relation ``$n$ evenly divides $m$'' is a
primitive recursive relation.

Primitive recursive functions---those that can be defined using just
primitive recursion---are not, however, the only computable
number-theoretic functions. Many generalizations of primitive
recursion have been considered, but the most powerful and
widely-accepted additional way of computing functions is by unbounded
search.  This leads to the definition of \emph{partial recursive
  functions}, and a related definition to \emph{general recursive
  functions}.  General recursive functions are computable and total,
and the definition characterizes exactly the partial recursive
functions that happen to be total.  Recursive functions can simulate
every other model of computation (Turing machines, lambda calculus,
etc.) and so represent one of the many accepted models of computation.



% ==== OLP file: content/computability/recursive-functions/non-pr-functions.tex ====
% Part: computability
% Chapter: recursive-functions
% Section: non-pr-functions



\olfileid{cmp}{rec}{npr}
\olsection{Non-Primitive Recursive Functions}

The primitive recursive functions do not exhaust the intuitively
computable functions. It should be intuitively clear that we can make
a list of all the unary primitive recursive functions,
$f_0$, $f_1$, $f_2$,~\dots such that we can effectively compute the value of
$f_x$ on input $y$; in other words, the function $g(x,y)$, defined by
\[
g(x,y) = f_x(y)
\]
is computable. But then so is the function
\begin{eqnarray*}
h(x) & = & g(x,x) + 1 \\
& = & f_x(x) +1.
\end{eqnarray*}
For each primitive recursive function $f_i$, the value of $h$ and
$f_i$ differ at $i$. So $h$ is computable, but not primitive
recursive; and one can say the same about $g$. This is an
``effective'' version of Cantor's diagonalization argument.

One can provide more explicit examples of computable functions that
are not primitive recursive. For example, let the notation $g^n(x)$
denote $g(g(\dots g(x)))$, with $n$ $g$'s in all; and define a
sequence $g_0,g_1,\dots$ of functions by
\begin{eqnarray*}
g_0(x) & = & x+1 \\
g_{n + 1}(x) & = & g_n^x(x)
\end{eqnarray*}
You can confirm that each function $g_n$ is primitive recursive. Each
successive function grows much faster than the one before; $g_1(x)$ is
equal to $2x$, $g_2(x)$ is equal to $2^x \cdot x$, and $g_3(x)$ grows
roughly like an exponential stack of $x$ $2$'s. The Ackermann--P\'eter
function is essentially the function $G(x) = g_x(x)$, and one can show
that this grows faster than any primitive recursive function.

Let us return to the issue of enumerating the primitive recursive
functions. Remember that we have assigned symbolic notations to each
primitive recursive function; so it suffices to enumerate
notations. We can assign a natural number $\#(F)$ to each notation $F$,
recursively, as follows:
\begin{eqnarray*}
\#(0) & = & \langle 0 \rangle \\
\#(S) & = & \langle 1 \rangle \\
\#(\Proj{n}{i}) & = & \langle 2, n, i \rangle \\
\#(\fn{Comp}_{k,l}[H,G_0,\dots,G_{k-1}]) & = & \langle
3,k,l,\#(H),\#(G_0),\dots,\#(G_{k-1}) \rangle \\
\#(\fn{Rec}_l[G,H]) & = & \langle 4, l, \#(G), \#(H) \rangle
\end{eqnarray*}
Here we are using the fact that every sequence of numbers can be viewed
as a natural number, using the codes from the last section. The upshot
is that every code is assigned a natural number. Of course, some
sequences (and hence some numbers) do not correspond to notations; but
we can let $f_i$ be the unary primitive recursive function with
notation coded as $i$, if $i$ codes such a notation; and the constant
$0$ function otherwise. The net result is that we have an explicit way of
enumerating the unary primitive recursive functions.

(In fact, some functions, like the constant zero function, will appear
more than once on the list. This is not just an artifact of our
coding, but also a result of the fact that the constant zero function has
more than one notation. We will later see that one can not computably
avoid these repetitions; for example, there is no computable function
that decides whether or not a given notation represents the constant
zero function.)

We can now take the function $g(x,y)$ to be given by $f_x(y)$, where
$f_x$ refers to the enumeration we have just described. How do we know
that $g(x,y)$ is computable? Intuitively, this is clear: to compute
$g(x,y)$, first ``unpack'' $x$, and see if it is a notation for a unary
function. If it is, compute the value of that function on input~$y$.

\begin{tagblock}{TMs}
\begin{digress}
You may already be convinced that (with some work!) one can write a
program (say, in Java or C++) that does this; and now we can appeal to
the Church--Turing thesis, which says that anything that, intuitively,
is computable can be computed by a Turing machine.

Of course, a more direct way to show that $g(x,y)$ is computable is to
describe a Turing machine that computes it, explicitly. This would,
in particular, avoid the Church--Turing thesis and appeals to
intuition. Soon we will have built up enough machinery to show
that $g(x,y)$ is computable, appealing to a model of computation that
can be \emph{simulated} on a Turing machine: namely, the recursive
functions.
\end{digress}
\end{tagblock}


% ==== OLP file: content/computability/recursive-functions/normal-form.tex ====
% Part: computability
% Chapter: recursive-functions
% Section: normal-form



\olfileid{cmp}{rec}{nft}
\olsection{The Normal Form Theorem}

\begin{thm}[Kleene's Normal Form Theorem]
\ollabel{thm:kleene-nf}
There is a primitive recursive relation $T(e, x, s)$ and a primitive
recursive function $U(s)$, with the following property: if $f$ is any
partial recursive function, then for some~$e$,
\[
f(x) \simeq U(\umin{s}{T(e, x, s)})
\]
for every $x$.
\end{thm}

\begin{explain}
The proof of the normal form theorem is involved, but the basic idea
is simple.  Every partial recursive function has an \emph{index}~$e$,
intuitively, a number coding its program or definition.  If $f(x)
\fdefined$, the computation can be recorded systematically and coded
by some number~$s$, and the fact that $s$ codes the computation of~$f$
on input~$x$ can be checked primitive recursively using only $x$ and
the definition~$e$.  Consequently, the relation~$T$, ``the function
with index~$e$ has a computation for input~$x$, and $s$ codes this
computation,'' is primitive recursive. Given the full record of the
computation~$s$, the ``upshot'' of~$s$ is the value of~$f(x)$, and it
can be obtained from~$s$ primitive recursively as well.

The normal form theorem shows that only a single unbounded search is
required for the definition of any partial recursive function.
Basically, we can search through all numbers until we find one that
codes a computation of the function with index~$e$ for input~$x$. We
can use the numbers~$e$ as ``names'' of partial recursive functions,
and write $\cfind{e}$ for the function~$f$ defined by the equation in
the theorem.  Note that any partial recursive function can have more
than one index---in fact, every partial recursive function has
infinitely many indices.
\end{explain}


% ==== OLP file: content/computability/recursive-functions/notation-pr-functions.tex ====
% Part: computability
% Chapter: recursive-functions
% Section: notations-pr-functions



\olfileid{cmp}{rec}{not}
\olsection{Primitive Recursion Notations}

One advantage to having the precise inductive description of the primitive
recursive functions is that we can be systematic in describing them.
For example, we can assign a ``notation'' to each such function, as
follows. Use symbols $\Zero$, $\Succ$, and $\Proj{n}{i}$ for zero,
successor, and the projections. Now suppose $h$ is defined by
composition from a $k$-place function~$f$ and $n$-place functions $g_0$,
\dots,~$g_{k-1}$, and we have assigned notations $F$, $G_0$,
\dots,~$G_{k-1}$ to the latter functions. Then, using a new symbol
$\fn{Comp}_{k,n}$, we can denote the function $h$ by
$\fn{Comp}_{k,n}[F,G_0,\dots,G_{k-1}]$. 

For functions defined by primitive recursion, we can use analogous
notations. Suppose the $(k+1)$-ary function~$h$ is defined by
primitive recursion from the $k$-ary function~$f$ and the $(k+2)$-ary
function~$g$, and the notations assigned to $f$ and~$g$ are $F$
and~$G$, respectively. Then the notation assigned to~$h$ is
$\fn{Rec}_k[F,G]$. 

Recall that the addition function is defined by primitive recursion as
\begin{align*}
  \Add(x_0, 0) & = \Proj{1}{0}(x_0) = x_0\\
  \Add(x_0, y+1) & = \Succ(\Proj{3}{2}(x_0, y, \Add(x_0, y))) = \Add(x_0, y) +1
\end{align*}
Here the role of~$f$ is played by $\Proj{1}{0}$, and the role of~$g$
is played by $\Succ(\Proj{3}{2}(x_0, y, z))$, which is assigned the
notation $\fn{Comp}_{1,3}[\Succ,\Proj{3}{2}]$ as it is the result of
defining a function by composition from the $1$-ary function~$\Succ$
and the $3$-ary function~$\Proj{3}{2}$. With this setup, we can denote
the addition function by
\[
\fn{Rec}_1[\Proj{1}{0},\fn{Comp}_{1,3}[\Succ,\Proj{3}{2}]].
\]
Having these notations sometimes proves useful, e.g., when enumerating
primitive recursive functions.

\begin{prob}
  Give the complete primitive recursive notation for~$\Mult$.
\end{prob}


% ==== OLP file: content/computability/recursive-functions/other-recursions.tex ====
% Part: computability
% Chapter: recursive-functions
% Section: other-recursions



\olfileid{cmp}{rec}{ore}
\olsection{Other Recursions}

Using pairing and sequencing, we can justify more exotic (and
useful) forms of primitive recursion. For example, it is often useful
to define two functions simultaneously, such as in the following
definition:
\begin{align*}
h_0(\vec x, 0) & = f_0(\vec x) \\
h_1(\vec x, 0) & = f_1(\vec x) \\
h_0(\vec x, y+1) & = g_0(\vec x, y, h_0(\vec x, y), h_1(\vec x, y)) \\
h_1(\vec x, y+1) & = g_1(\vec x, y, h_0(\vec x, y), h_1(\vec x, y))
\end{align*}
This is an instance of \emph{simultaneous recursion}. Another useful
way of defining functions is to give the value of $h(\vec x, y+1)$ in
terms of \emph{all} the values $h(\vec x, 0)$, \dots,~$h(\vec x, y)$, as in
the following definition:
\begin{align*}
h(\vec x, 0) & = f(\vec x) \\
h(\vec x, y+1) & = g(\vec x, y, \tuple{h(\vec x, 0), \dots, h(\vec x, y)}).
\end{align*}
The following schema captures this idea more succinctly:
\[
h(\vec x, y) = g(\vec x, y, \tuple{h(\vec x, 0), \dots, h(\vec x, y-1)})
\]
with the understanding that the last argument to $g$ is just the
empty sequence when $y$ is $0$. In either formulation, the idea is
that in computing the ``successor step,'' the function $h$ can make
use of the entire sequence of values computed so far.
This is known as a \emph{course-of-values} recursion. For a particular
example, it can be used to justify the following type of definition:
\begin{align*}
h(\vec x, y) & = \begin{cases}
  g(\vec x, y, h(\vec x, k(\vec x, y))) & \text{if $k(\vec x, y) < y$} \\
  f(\vec x) & \text{otherwise}
\end{cases}
\end{align*}
In other words, the value of $h$ at $y$ can be computed in terms of
the value of $h$ at \emph{any} previous value, given by~$k$.


\begin{prob}
  Define the remainder function $r(x,y)$ by course-of-values
  recursion. (If $x$, $y$ are natural numbers and $y > 0$, $r(x,y)$ is
  the number less than~$y$ such that $x = z\times y + r(x,y)$ for
  some~$z$. For definiteness, let's say that if $y=0$, $r(x,0) = 0$.)
\end{prob}

You should think about how to obtain these functions using ordinary
primitive recursion. One final version of primitive recursion is more
flexible in that one is allowed to change the \emph{parameters} (side
values) along the way:
\begin{align*}
h(\vec x, 0) & = f(\vec x) \\
h(\vec x, y+1) & = g(\vec x, y, h(k(\vec x), y))
\end{align*}
This, too, can be simulated with ordinary primitive recursion. (Doing
so is tricky. For a hint, try unwinding the computation by hand.)


% ==== OLP file: content/computability/recursive-functions/partial-functions.tex ====
% Part: computability
% Chapter: recursive-functions
% Section: partial-functions



\olfileid{cmp}{rec}{par}
\olsection{Partial Recursive Functions}


To motivate the definition of the recursive functions, note that our
proof that there are computable functions that are not primitive
recursive actually establishes much more. The argument was
simple: all we used was the fact that it is possible to enumerate
functions $f_0,f_1,\dots$ such that, as a function of $x$ and $y$,
$f_x(y)$ is computable. So the argument applies to \emph{any class of
  functions that can be enumerated in such a way}. This puts us in a
bind: we would like to describe the computable functions explicitly;
but any explicit description of a collection of computable functions
cannot be exhaustive!

The way out is to allow \emph{partial} functions to come into play. We
will see that it \emph{is} possible to enumerate the partial
computable functions.\iftag{TMs}{ In fact, we already pretty much know
  that this is the case, since it is possible to enumerate Turing
  machines in a systematic way.}{} We will come back to our diagonal
argument later, and explore why it does not go through when partial
functions are included.

The question is now this: what do we need to add to the primitive
recursive functions to obtain all the partial recursive functions? We
need to do two things:
\begin{enumerate}
\item Modify our definition of the primitive recursive functions to
  allow for partial functions as well.
\item \emph{Add} something to the definition, so that some new partial
  functions are included.
\end{enumerate}

The first is easy. As before, we will start with zero, successor, and
projections, and close under composition and primitive recursion. The
only difference is that we have to modify the definitions of
composition and primitive recursion to allow for the possibility that
some of the terms in the definition are not defined. If $f$ and $g$
are partial functions, we will write $f(x) \fdefined$ to mean that $f$
is defined at $x$, i.e., $x$ is in the domain of $f$; and $f(x)
\fundefined$ to mean the opposite, i.e., that $f$ is not defined at~$x$.
We will use $f(x) \simeq g(x)$ to mean that either $f(x)$ and $g(x)$
are both undefined, or they are both defined and equal. We will use these
notations for more complicated terms as well. We will adopt the
convention that if $h$ and $g_0$, \dots,~$g_k$ all are partial functions,
then
\[
h(g_0(\vec x),\dots,g_k(\vec x))
\]
is defined if and only if each $g_i$ is defined at $\vec x$, and $h$
is defined at $g_0(\vec x)$, \dots,~$g_k(\vec x)$. With this
understanding, the definitions of composition and primitive recursion
for partial functions is just as above, except that we have to replace
``$=$'' by ``$\simeq$''.

What we will add to the definition of the primitive recursive
functions to obtain partial functions is the \emph{unbounded search
  operator}. If $f(x,\vec z)$ is any partial function on the natural
numbers, define $\mu x \; f(x,\vec z)$ to be
\begin{quote}
  the least $x$ such that $f(0,\vec z), f(1,\vec z), \dots, f(x,\vec
  z)$ are all defined, and $f(x,\vec z) = 0$, if such an $x$ exists
\end{quote}
with the understanding that $\mu x \; f(x,\vec z)$ is undefined
otherwise. This defines $\mu x \; f(x,\vec z)$ uniquely.

\begin{explain}
Note that our definition makes no reference to\iftag{TMs}{ Turing
  machines, or}{} algorithms, or any specific computational model. But
like composition and primitive recursion, there is an operational,
computational intuition behind unbounded search. When it
comes to the computability of a partial function, arguments
where the function is undefined correspond to inputs for which the
computation does not halt. The procedure for computing $\mu x \;
f(x,\vec z)$ will amount to this: compute $f(0,\vec z), f(1,\vec z),
f(2,\vec z)$ until a value of 0 is returned. If any of the
intermediate computations do not halt, however, neither does the
computation of $\mu x \; f(x,\vec z)$.
\end{explain}

If $R(x,\vec z)$ is any relation, $\mu x \; R(x,\vec z)$ is defined to
be $\mu x \; (1 \tsub \Char{R}(x,\vec z))$. In other words, $\mu x \;
R(x,\vec z)$ returns the least value of $x$ such that $R(x,\vec z)$
holds. So, if $f(x,\vec z)$ is a total function, $\mu x \; f(x,\vec
z)$ is the same as $\mu x \; (f(x,\vec z) = 0)$. But note that our
original definition is more general, since it allows for the
possibility that $f(x,\vec z)$ is not everywhere defined (whereas, in
contrast, the characteristic function of a relation is always total).

\begin{defn}
The set of \emph{partial recursive functions} is the smallest set of
partial functions from the natural numbers to the natural numbers (of
various arities) containing zero, successor, and projections, and
closed under composition, primitive recursion, and unbounded search.
\end{defn}

Of course, some of the partial recursive functions will happen to be
total, i.e., defined for every argument.

\begin{defn}
\ollabel{defn:recursive-fn}
The set of \emph{recursive functions} is
the set of partial recursive functions that are total.
\end{defn}

A recursive function is sometimes called ``total recursive'' to
emphasize that it is defined everywhere.




% ==== OLP file: content/computability/recursive-functions/pr-functions-computable.tex ====
% Part: computability
% Chapter: recursive-functions
% Section: pr-functions-computable



\olfileid{cmp}{rec}{cmp}
\olsection{Primitive Recursive Functions are Computable}

Suppose a function $h$ is defined by primitive recursion
\begin{eqnarray*}
h(\vec x, 0)     & = & f(\vec x) \\
h(\vec x, y + 1) & = & g(\vec x, y, h(\vec x, y))
\end{eqnarray*}
and suppose the functions $f$ and $g$ are computable.  (We use $\vec x$ to abbreviate $x_0$, \dots, $x_{k-1}$.) Then $h(\vec
x, 0)$ can obviously be computed, since it is just $f(\vec x)$ which we
assume is computable.  $h(\vec x, 1)$ can then also be computed, since
$1 = 0 + 1$ and so $h(\vec x, 1)$ is just
\begin{align*}
 h(\vec x, 1) & = g(\vec x, 0, h(\vec x, 0)) =  g(\vec x, 0, f(\vec x)).
\intertext{We can go on in this way and  compute}
h(\vec x, 2) & = g(\vec x, 1, h(\vec x, 1)) = g(\vec x, 1, g(\vec x, 0, f(\vec x)))\\
h(\vec x, 3) & = g(\vec x, 2, h(\vec x, 2)) = g(\vec x, 2, g(\vec x, 1, g(\vec x, 0, f(\vec x))))\\
h(\vec x, 4) & = g(\vec x, 3, h(\vec x, 3)) = g(\vec x, 3, g(\vec x, 2, g(\vec x, 1, g(\vec x, 0, f(\vec x)))))\\
& \vdots
\end{align*}
Thus, to compute $h(\vec x, y)$ in general, successively compute
$h(\vec x, 0)$, $h(\vec x, 1)$, \dots, until we reach $h(\vec x, y)$.

Thus, a primitive recursive definition yields a new computable
function if the functions $f$ and $g$ are computable.  Composition of
functions also results in a computable function if the functions $f$
and $g_i$ are computable.

Since the basic functions $\Zero$, $\Succ$, and $\Proj{n}{i}$ are
computable, and composition and primitive recursion yield computable
functions from computable functions, this means that every primitive
recursive function is computable.


% ==== OLP file: content/computability/recursive-functions/pr-functions.tex ====
% Part: computability
% Chapter: recursive-functions
% Section: pr-functions



\olfileid{cmp}{rec}{prf}
\olsection{Primitive Recursion Functions}

Let us record again how we can define new functions from existing ones
using primitive recursion and composition.

\begin{defn}
  \ollabel{defn:primitive-recursion} Suppose $f$ is a $k$-place
  function ($k\ge 1$) and $g$ is a $(k+2)$-place function. The
  function defined by \emph{primitive recursion from $f$ and $g$} is
  the $(k+1)$-place function~$h$ defined by the equations
  \begin{align*}
    h(x_0,\dots,x_{k-1},0) & =  f(x_0,\dots,x_{k-1}) \\
    h(x_0,\dots,x_{k-1},y+1) & = g(x_0,\dots,x_{k-1}, y, h(x_0,\dots,x_{k-1}, y))
  \end{align*}
\end{defn}

\begin{defn}
  \ollabel{defn:composition} Suppose $f$ is a $k$-place function, and
  $g_0$, \dots, $g_{k-1}$ are $k$ functions which are all
  $n$-place. The function defined by \emph{composition from $f$ and
    $g_0$, \dots,~$g_{k-1}$} is the $n$-place function~$h$ defined by
  \[
  h(x_0, \dots, x_{n-1}) =
  f(g_0(x_0, \dots, x_{n-1}), \dots, g_{k-1}(x_0, \dots, x_{n-1})).
  \]
\end{defn}

In addition to $\Succ$ and the projection functions
\[
\Proj{n}{i}(x_0,\dots,x_{n-1}) = x_i,
\]
for each natural number $n$ and $i < n$, we will include among the
primitive recursive functions the function~$\Zero(x) = 0$.

\begin{defn}
  The set of primitive recursive functions is the set of functions 
  from $\Nat^n$ to $\Nat$, defined inductively by the
  following clauses:
  \begin{enumerate}
  \item $\Zero$ is primitive recursive.
  \item $\Succ$ is primitive recursive.
  \item Each projection function $\Proj{n}{i}$ is primitive recursive.
  \item If $f$ is a $k$-place primitive recursive function and $g_0$,
    \dots,~$g_{k-1}$ are $n$-place primitive recursive functions, then
    the composition of $f$ with $g_0$, \dots,~$g_{k-1}$ is primitive
    recursive.
  \item If $f$ is a $k$-place primitive recursive function and $g$ is
    a $k+2$-place primitive recursive function, then the function
    defined by primitive recursion from $f$ and $g$ is primitive
    recursive.
\end{enumerate}
\end{defn}

\begin{explain}
Put more concisely, the set of primitive recursive functions is the
smallest set containing $\Zero$, $\Succ$, and the projection
functions~$\Proj{n}{j}$, and which is closed under composition and
primitive recursion.

Another way of describing the set of primitive recursive functions is
by defining it in terms of ``stages.'' Let $S_0$ denote the set of
starting functions: $\Zero$, $\Succ$, and the projections. These are
the primitive recursive functions of stage~$0$. Once a stage $S_i$ has
been defined, let $S_{i+1}$ be the set of all functions you get by
applying a single instance of composition or primitive recursion to
functions already in~$S_i$. Then
\[
S = \bigcup_{i \in \Nat} S_i
\]
is the set of all primitive recursive functions
\end{explain}

Let us verify that $\Add$ is a primitive recursive function.

\begin{prop}
  The addition function $\Add(x,y) = x+y$ is primitive recursive.
\end{prop}

\begin{proof}
We already have a primitive recursive definition of $\Add$ in terms of
two functions $f$ and~$g$ which matches the format of
\olref{defn:primitive-recursion}:
\begin{align*}
  \Add(x_0, 0) & = f(x_0) = x_0\\
  \Add(x_0, y+1) & = g(x_0, y, \Add(x_0, y)) =
  \Succ(\Add(x_0, y))
\end{align*}
So $\Add$ is primitive recursive provided $f$ and $g$ are as
well. $f(x_0) = x_0 = \Proj{1}{0}(x_0)$, and the projection functions
count as primitive recursive, so $f$ is primitive recursive. The
function $g$ is the three-place function $g(x_0, y, z)$ defined by
\[
g(x_0, y, z) = \Succ(z).
\]
This does not yet tell us that $g$ is primitive recursive, since $g$
and $\Succ$ are not quite the same function: $\Succ$ is one-place, and
$g$ has to be three-place. But we can define $g$ ``officially'' by
composition as
\[
g(x_0, y, z) = \Succ(\Proj{3}{2}(x_0, y, z))
\]
Since $\Succ$ and $\Proj{3}{2}$ count as primitive recursive
functions, $g$ does as well, since it can be defined by composition
from primitive recursive functions.
\end{proof}

\begin{prop}
  \ollabel{prop:mult-pr}
  The multiplication function $\Mult(x,y) = x \cdot y$ is primitive recursive.
\end{prop}

\begin{proof}
  Exercise.
\end{proof}

\begin{prob}
  Prove \olref[cmp][rec][prf]{prop:mult-pr} by showing that the
  primitive recursive definition of $\Mult$ can be put into the
  form required by \olref[cmp][rec][prf]{defn:primitive-recursion} and
  showing that the corresponding functions $f$ and~$g$ are primitive
  recursive.
\end{prob}

\begin{ex}
Here's our very first example of a primitive recursive definition:
\begin{align*}
h(0) & =  1 \\
h(y+1) & =  2 \cdot h(y).
\end{align*}
This function cannot fit into the form required by
\olref{defn:primitive-recursion}, since $k=0$. The definition also
involves the constants $1$ and $2$. To get around the first problem,
let's introduce a dummy argument and define the function~$h'$:
\begin{align*}
h'(x_0, 0) & =  f(x_0) = 1 \\
h'(x_0, y+1) & =  g(x_0, y, h'(x_0, y)) = 2 \cdot h'(x_0, y).
\end{align*}
The function $f(x_0) = 1$ can be defined from $\Succ$ and $\Zero$ by
composition: $f(x_0) = \Succ(\Zero(x_0))$. The function $g$ can be
defined by composition from $g'(z) = 2 \cdot z$ and projections:
\begin{align*}
g(x_0, y, z) & = g'(\Proj{3}{2}(x_0, y, z))
\intertext{and $g'$ in turn can be defined by composition as}
g'(z) & = \Mult(g''(z), \Proj{1}{0}(z))
\intertext{and}
g''(z) & = \Succ(f(z)),
\end{align*}
where $f$ is as above: $f(z) = \Succ(\Zero(z))$. Now that we
have~$h'$, we can use composition again to let $h(y) =
h'(\Proj{1}{0}(y),\Proj{1}{0}(y))$. This shows that $h$ can be defined
from the basic functions using a sequence of compositions and
primitive recursions, so $h$~is primitive recursive.
\end{ex}


% ==== OLP file: content/computability/recursive-functions/pr-relations.tex ====
% Part: computability
% Chapter: recursive-functions
% Section: pr-relations



\olfileid{cmp}{rec}{prr}
\olsection{Primitive Recursive Relations}


\begin{defn}
A relation $R(\vec x)$ is said to be primitive recursive if its characteristic
function,
\[
\Char{R}(\vec x) = \left\{
  \begin{array}{ll}
  1 & \mbox{if $R(\vec x)$} \\
  0 & \mbox{otherwise}
  \end{array}
\right.
\]
is primitive recursive.
\end{defn}

In other words, when one speaks of a primitive recursive relation
$R(\vec x)$, one is referring to a relation of the form $\Char{R}(\vec
x) = 1$, where $\Char{R}$ is a primitive recursive function which, on
any input, returns either 1 or 0. For example, the relation
$\fn{IsZero}(x)$, which holds if and only if $x = 0$, corresponds to the
function $\Char{\fn{IsZero}}$, defined using primitive recursion by
\begin{align*}
\Char{\fn{IsZero}}(0) & = 1,\\
\Char{\fn{IsZero}}(x+1) & = 0.
\end{align*}

It should be clear that one can compose relations with other primitive
recursive functions. So the following are also primitive recursive:
\begin{enumerate}
\item The equality relation, $x = y$, defined by $\fn{IsZero}(\left|x -
  y\right|)$
\item The less-than relation, $x \leq y$, defined by $\fn{IsZero}(x
  \tsub y)$
\end{enumerate}


\begin{prop}
  The set of primitive recursive relations is closed under Boolean
  operations, that is, 
  if $P(\vec x)$ and $Q(\vec x)$ are primitive recursive, so are
  \begin{enumerate}
  \item $\lnot P(\vec x)$
  \item $P(\vec x) \land Q(\vec x)$
  \item $P(\vec x) \lor Q(\vec x)$
  \item $P(\vec x) \lif Q(\vec x)$
  \end{enumerate}
\end{prop}

\begin{proof}
  Suppose $P(\vec x)$ and $Q(\vec x)$ are primitive recursive, i.e.,
  their characteristic functions $\Char{P}$ and $\Char{Q}$ are.  We
  have to show that the characteristic functions of $\lnot P(\vec x)$,
  etc., are also primitive recursive.
  \[
  \Char{\lnot P}(\vec x) = \begin{cases}
    0 & \text{if $\Char{P}(\vec x) = 1$}\\
    1 & \text{otherwise}
  \end{cases}
  \]
  We can define $\Char{\lnot P}(\vec x)$ as $1 \tsub \Char{P}(\vec x)$.
  \[
  \Char{P \land Q}(\vec x) = \begin{cases}
    1 & \text{if $\Char{P}(\vec x) = \Char{Q}(\vec x) = 1$}\\
    0 & \text{otherwise}
  \end{cases}
  \]
  We can define $\Char{P \land Q}(\vec x)$ as $\Char{P}(\vec x) \cdot
  \Char{Q}(\vec x)$ or as $\fn{min}(\Char{P}(\vec x), \Char{Q}(\vec
  x))$. Similarly, 
  \begin{align*}
    \Char{P \lor Q}(\vec x) & = \fn{max}(\Char{P}(\vec x), \Char{Q}(\vec x)) \text{ and}\\
    \Char{P \lif Q}(\vec x) & = \fn{max}(1 \tsub
  \Char{P}(\vec x), \Char{Q}(\vec x)).
  \end{align*}
\end{proof}

\begin{prop}
  The set of primitive recursive relations is closed under bounded
  quantification, i.e., if $R(\vec x, z)$ is a primitive recursive
  relation, then so are the relations
  \begin{align*}
    & \bforall{z < y}{R(\vec x, z)} \text{ and}\\
    & \bexists{z < y}{R(\vec x, z)}.
  \end{align*}
  $\bforall{z < y}{R(\vec x, z)}$ holds of $\vec x$ and $y$ if and
  only if $R(\vec x, z)$ holds for every~$z$ less than~$y$, and
  similarly for $\bexists{z < y}{R(\vec x, z)}$.
\end{prop}

\begin{proof}
  By convention, we take $\bforall{z < 0}{R(\vec x, z)}$ to be true
  (for the trivial reason that there are no $z$ less than~$0$) and
  $\bexists{z < 0}{R(\vec x, z)}$ to be false. A bounded
  universal quantifier functions just like a finite product or
  iterated minimum, i.e., if $P(\vec x, y) \defiff \bforall{z <
  y}{R(\vec x, z)}$ then $\Char{P}(\vec x, y)$ can be defined by
  \begin{align*}
    \Char{P}(\vec x, 0) & = 1\\
    \Char{P}(\vec x, y+1) & =
    \fn{min}(\Char{P}(\vec x, y), \Char{R}(\vec x, y)).
  \end{align*}
  Bounded existential quantification can similarly be defined using
  $\fn{max}$. Alternatively, it can be defined from bounded universal
  quantification, using the equivalence $\bexists{z < y}{R(\vec x, z)}
  \liff \lnot \bforall{z < y}{\lnot R(\vec x, z)}$. Note that, for
  example, a bounded quantifier of the form $\bexists{x \leq y}{\dots
  x\dots}$ is equivalent to $\bexists{x < y+1}{\dots x \dots}$.
\end{proof}

\begin{prob}
  Show that the three place relation $x \equiv y \mod n$ (congruence
  modulo~$n$) is primitive recursive.
\end{prob}

Another useful primitive recursive function is the conditional
function, $\fn{cond}(x,y,z)$, defined by
\begin{align*}
  \fn{cond}(x,y,z) & = \begin{cases}
  y & \text{if $x = 0$} \\
  z & \text{otherwise}.
\end{cases}
\intertext{This is defined recursively by}
\fn{cond}(0,y,z) & = y,\\
\fn{cond}(x+1,y,z) & = z.
\end{align*}
One can use this to justify definitions of primitive recursive functions
by cases from primitive recursive relations:

\begin{prop}
If $g_0(\vec x)$, \dots,~$g_m(\vec x)$ are primitive recursive functions, and $R_0(\vec
x)$, \dots, $R_{m-1}(\vec x)$ are primitive recursive relations, then
the function $f$ defined by
\[
f(\vec x) = \begin{cases}
    g_0(\vec x) & \text{if $R_0(\vec{x})$} \\
    g_1(\vec x) & \text{if $R_1(\vec{x})$ and not $R_0(\vec{x})$} \\
    \vdots & \\
    g_{m-1}(\vec x) & \text{if $R_{m-1}(\vec{x})$ and none of the
      previous hold}
    \\
    g_m(\vec x) & \mbox{otherwise}
\end{cases}
\]
is also primitive recursive.
\end{prop}

\begin{proof}
  When $m = 1$, this is just the function defined by
  \[
  f(\vec x) = \fn{cond}(\Char{\lnot R_0}(\vec x),g_0(\vec x),g_1(\vec
  x)).
  \]
  For $m$ greater than $1$, one can just compose definitions of this
  form.
\end{proof}

% ==== OLP file: content/computability/recursive-functions/primes.tex ====
% Part: computability
% Chapter: recursive-functions
% Section: primes



\olfileid{cmp}{rec}{pri}
\olsection{Primes}

Bounded quantification and bounded minimization provide us with a good
deal of machinery to show that natural functions and relations are
primitive recursive. For example, consider the relation ``$x$
divides $y$'', written $x \mid y$.  The relation $x \mid y$ holds if
division of $y$ by~$x$ is possible without remainder, i.e., if $y$ is
an integer multiple of~$x$.  (If it doesn't hold, i.e., the remainder
when dividing $x$ by $y$ is $> 0$, we write $x \nmid y$.) In other
words, $x \mid y$ iff for some~$z$, $x \cdot z = y$.  Obviously, any
such $z$, if it exists, must be $\leq y$. So, we have that $x \mid y$
iff for some $z \le y$, $x \cdot z = y$.  We can define the relation
$x \mid y$ by bounded existential quantification from $=$ and
multiplication by
\[
x \mid y \defiff \bexists{z \leq y}{(x \cdot z) = y}.
\]
We've thus shown that $x \mid y$ is primitive recursive.

A natural number~$x$ is \emph{prime} if it is neither $0$ nor $1$ and
is only divisible by $1$ and itself. In other words, prime numbers are
such that, whenever $y \mid x$, either $y = 1$ or~$y=x$.  To test if
$x$~is prime, we only have to check if $y \mid x$ for all $y \le x$,
since if $y > x$, then automatically~$y \nmid x$.  So, the relation
$\fn{Prime}(x)$, which holds iff $x$ is prime, can be defined by
\[
\fn{Prime}(x) \defiff x \geq 2 \land \bforall{y \leq x}{(y \mid x \lif y
  = 1 \lor y = x)}
\]
and is thus primitive recursive.

The primes are $2$, $3$, $5$, $7$, $11$, etc. Consider the function
$p(x)$ which returns the $x$th prime in that sequence, i.e., $p(0) =
2$, $p(1) = 3$, $p(2) = 5$, etc. (For convenience we will often write
$p(x)$ as $p_x$ ($p_0=2$, $p_1=3$, etc.)

If we had a function
$\fn{nextPrime(x)}$, which returns the first prime number larger
than~$x$, $p$~can be easily defined using primitive recursion:
\begin{align*}
  p(0) & = 2\\
  p(x+1) & = \fn{nextPrime}(p(x))
\end{align*}
Since $\fn{nextPrime}(x)$ is the least $y$ such that $y > x$ and
$y$~is prime, it can be easily computed by unbounded search. But it
can also be defined by bounded minimization, thanks to a result due to
Euclid: there is always a prime number between $x$ and $\fact{x}+1$.
\[
  \fn{nextPrime(x)} =
  \bmin{y \leq \fact{x}+1}{(y > x \land \fn{Prime}(y))}.
\]
This shows, that $\fn{nextPrime}(x)$ and hence $p(x)$ are (not just
computable but) primitive recursive.

(If you're curious, here's a quick proof of Euclid's theorem. Suppose
$p_n$ is the largest prime $\le x$ and consider the product $p =
p_0\cdot p_1 \cdot \dots \cdot p_n$ of all primes~$\le x$. Either
$p+1$ is prime or there is a prime between $x$ and~$p+1$.  Why?
Suppose $p+1$ is not prime. Then some prime number $q \mid p+1$ where
$q < p+1$. None of the primes $\le x$ divide $p+1$. (By definition
of~$p$, each of the primes $p_i \le x$ divides~$p$, i.e., with
remainder~$0$. So, each of the primes $p_i \le x$ divides $p+1$ with
remainder~$1$, and so $p_i \nmid p+1$.)  Hence, $q$ is a prime $>
x$ and $< p+1$.  And $p \le \fact{x}$, so there is a prime $> x$ and $\le
\fact{x}+1$.)

\begin{prob}
Define integer division $d(x, y)$ using bounded minimization.
\end{prob}


% ==== OLP file: content/computability/recursive-functions/primitive-recursion.tex ====
% Part: computability
% Chapter: recursive-functions
% Section: primitive-recursion



\olfileid{cmp}{rec}{pre}
\olsection{Primitive Recursion}

A characteristic of the natural numbers is that every natural number
can be reached from~$0$ by applying the successor operation~$+1$
finitely many times---any natural number is either~$0$ or the
successor of \dots{} the successor of~$0$. One way to specify a
function~$h\colon\Nat \to \Nat$ that makes use of this fact is this:
(a)~specify what the value of $h$ is for argument~$0$, and (b)~also
specify how to, given the value of $h(x)$, compute the value of
$h(x+1)$. For (a) tells us directly what $h(0)$ is, so $h$~is defined
for~$0$. Now, using the instruction given by (b) for $x=0$, we can
compute $h(1) = h(0+1)$ from~$h(0)$. Using the same instructions for
$x=1$, we compute $h(2) = h(1+1)$ from~$h(1)$, and so on. For every
natural number~$x$, we'll eventually reach the step where we define
$h(x)$ from $h(x+1)$, and so $h(x)$ is defined for all $x \in \Nat$.

For instance, suppose we specify $h\colon \Nat \to \Nat$ by the following
two equations:
\begin{align*}
h(0) & =  1\\
h(x+1) & =  2 \cdot h(x)
\end{align*}
If we already know how to multiply, then these equations give us the
information required for (a) and~(b) above. By successively applying
the second equation, we get that
\begin{align*}
  h(1) & = 2\cdot h(0) = 2,\\
  h(2) & = 2\cdot h(1) = 2\cdot 2,\\
  h(3) & = 2 \cdot h(2) = 2\cdot 2 \cdot 2,\\
  & \vdots
\end{align*}
We see that the function~$h$ we have specified is $h(x) = 2^x$.

The characteristic feature of the natural numbers guarantees that
there is only one function~$h$ that meets these two criteria.  A pair
of equations like these is called a \emph{definition by primitive
recursion} of the function~$h$. It is so-called because we define~$h$
``recursively,'' i.e., the definition, specifically the second
equation, involves $h$ itself on the right-hand-side. It is
``primitive'' because in defining $h(x+1)$ we only use the
value~$h(x)$, i.e., the immediately preceding value. This is the
simplest way of defining a function on~$\Nat$ recursively.

We can define even more fundamental functions like addition and
multiplication by primitive recursion. In these cases, however, the
functions in question are $2$-place. We fix one of the argument
places, and use the other for the recursion. E.g, to define
$\Add(x, y)$ we can fix~$x$ and define the value first for $y=0$
and then for $y+1$ in terms of~$y$. Since $x$ is fixed, it will appear
on the left and on the right side of the defining equations.
\begin{align*}
\Add(x,0) & =  x\\
\Add(x,y+1) & =  \Add(x,y)+1
\end{align*}
These equations specify the value of $\Add$ for all $x$
\emph{and}~$y$. To find $\Add(2,3)$, for instance, we apply the
defining equations for $x = 2$, using the first to find
$\Add(2,0) = 2$, then using the second to successively find
$\Add(2,1) = 2 + 1 = 3$, $\Add(2, 2) = 3 + 1 = 4$, $\Add(2,
3) = 4 + 1 = 5$.

In the definition of $\Add$ we used $+$ on the right-hand-side of
the second equation, but only to add~$1$. In other words, we used the
successor function $\Succ(z) = z+1$ and applied it to the previous value
$\Add(x,y)$ to define $\Add(x,y+1)$. So we can think of the
recursive definition as given in terms of a single function which we
apply to the previous value. However, it doesn't hurt---and sometimes
is necessary---to allow the function to depend not just on the previous
value but also on $x$ and~$y$. Consider:
\begin{align*}
  \Mult(x,0) & =  0 \\
  \Mult(x,y+1) & =  \Add(\Mult(x,y),x)
\end{align*}
This is a primitive recursive definition of a function $\Mult$ by
applying the function $\Add$ to both the preceding value
$\Mult(x,y)$ and the first argument~$x$. It also defines the
function~$\Mult(x,y)$ for all arguments $x$ and~$y$. For instance,
$\Mult(2,3)$ is determined by successively computing $\Mult(2,0)$,
$\Mult(2,1)$, $\Mult(2,2)$, and~$\Mult(2,3)$:
\begin{align*}
  \Mult(2,0) & = 0\\
  \Mult(2,1) & = \Mult(2,0+1) =
  \Add(\Mult(2,0), 2) = \Add(0, 2) = 2\\
  \Mult(2,2) & = \Mult(2,1+1) =
  \Add(\Mult(2,1), 2) = \Add(2, 2) = 4\\
  \Mult(2,3) & = \Mult(2,2+1) =
  \Add(\Mult(2,2), 2) = \Add(4, 2) = 6
\end{align*}

The general pattern then is this: to give a primitive recursive
definition of a function~$h(x_0, \dots, x_{k-1}, y)$, we provide two
equations. The first defines the value of $h(x_0, \dots, x_{k-1}, 0)$
without reference to~$h$. The second defines the value of $h(x_0,
\dots, x_{k-1}, y+1)$ in terms of $h(x_0, \dots, x_{k-1}, y)$, the
other arguments $x_0$, \dots,~$x_{k-1}$, and~$y$. Only the immediately
preceding value of~$h$ may be used in that second equation.  If we
think of the operations given by the right-hand-sides of these two
equations as themselves being functions $f$ and~$g$, then the general
pattern to define a new function~$h$ by primitive recursion is this:
\begin{align*}
  h(x_0, \dots, x_{k-1}, 0) & = f(x_0, \dots, x_{k-1})\\
  h(x_0, \dots, x_{k-1}, y+1) & =
  g(x_0, \dots, x_{k-1}, y, h(x_0, \dots, x_{k-1}, y))
\end{align*}  
In the case of $\Add$, we have $k=1$ and $f(x_0) = x_0$ (the
identity function), and $g(x_0, y, z) = z + 1$ (the $3$-place function
that returns the successor of its third argument):
\begin{align*}
  \Add(x_0, 0) & = f(x_0) = x_0\\
  \Add(x_0, y+1) & = g(x_0, y, \Add(x_0, y)) =
  \Succ(\Add(x_0, y))
\end{align*}
In the case of $\Mult$, we have $f(x_0) = 0$ (the constant
function always returning~$0$) and $g(x_0, y, z) = \Add(z,x_0)$
(the $3$-place function that returns the sum of its last and first
argument):
\begin{align*}
  \Mult(x_0,0) & =  f(x_0) = 0 \\
  \Mult(x_0,y+1) & =  g(x_0, y, \Mult(x_0,y)) =
  \Add(\Mult(x_0,y), x_0)
\end{align*}


% ==== OLP file: content/computability/recursive-functions/recursive-functions.tex ====
% Part: computability
% Chapter: recursive-functions



\olchapter{cmp}{rec}{Recursive Functions}

\begin{editorial}
  These are Jeremy Avigad's notes on recursive functions, revised and
  expanded by Richard Zach. This chapter does contain some exercises,
  and can be included independently to provide the basis for a
  discussion of arithmetization of syntax.
\end{editorial}



\olimport{introduction}

\olimport{primitive-recursion}

\olimport{composition}

\olimport{pr-functions}

\olimport{notation-pr-functions}

\olimport{pr-functions-computable}

\olimport{examples}

\olimport{pr-relations}

\olimport{bounded-minimization}

\olimport{primes}

\olimport{sequences}

\olimport{trees}

\olimport{other-recursions}

\olimport{non-pr-functions}

\olimport{partial-functions}

\olimport{normal-form}

\olimport{halting-problem}

\olimport{general-recursive-functions}

\OLEndChapterHook


% ==== OLP file: content/computability/recursive-functions/sequences.tex ====
% Part: computability
% Chapter: recursive-functions
% Section: sequences



\olfileid{cmp}{rec}{seq}
\olsection{Sequences}

The set of primitive recursive functions is remarkably robust. But we
will be able to do even more once we have developed a adequate means
of handling \emph{sequences}. We will identify finite sequences of
natural numbers with natural numbers in the following way: the
sequence $\langle a_0, a_1, a_2, \dots, a_k \rangle$ corresponds to
the number
\[
p_0^{a_0+1} \cdot p_1^{a_1+1} \cdot p_2^{a_2+1} \cdot \dots \cdot
p_k^{a_k+1}.
\]
We add one to the exponents to guarantee that, for example, the
sequences $\langle 2, 7, 3\rangle$ and $\langle 2, 7, 3, 0, 0 \rangle$
have distinct numeric codes. We can take both $0$ and~$1$ to code the
empty sequence; for concreteness, let $\emptyseq$ denote~$0$.

The reason that this coding of sequences works is the so-called
Fundamental Theorem of Arithmetic: every natural number $n \ge 2$ can
be written in one and only one way in the form
\[
n = p_0^{a_0} \cdot p_1^{a_1} \cdot \dots \cdot p_k^{a_k}
\]
with $a_k \ge 1$. This guarantees that the mapping $\tuple{}(a_0,
\dots, a_k) = \tuple{a_0, \dots, a_k}$ is injective: different
sequences are mapped to different numbers; to each number only at most
one sequence corresponds.

We'll now show that the operations of determining the length of a
sequence, determining its $i$th element, appending an element to a
sequence, and concatenating two sequences, are all primitive
recursive.

\begin{prop}
  The function $\len{s}$, which returns the length of the sequence
  $s$, is primitive recursive.
\end{prop}

\begin{proof}
  Let $R(i, s)$ be the relation defined by
  \[
  R(i, s) \text{ iff }
  p_i \mid s \land p_{i+1} \nmid s.
  \]
  $R$ is clearly primitive recursive. Whenever $s$ is the code of a
  non-empty sequence, i.e.,
  \[
  s = p_0^{a_0+1} \cdot \dots \cdot p_{k}^{a_{k}+1},
  \]
  $R(i,s)$ holds if $p_i$ is the largest prime such that
  $p_i \mid s$, i.e., $i = k$. The length of $s$ thus is $i+1$ iff
  $p_i$ is the largest prime that divides~$s$, so we can let
  \[
  \len{s} =
  \begin{cases}
    0 & \text{if $s = 0$ or $s = 1$} \\
    1 + \bmin{i < s}{R(i, s)} & \text{otherwise}
  \end{cases}
  \]
  We can use bounded minimization here, since there is only one $i$ that
  satisfies $R(i,s)$ when $s$~is a code of a sequence, and if $i$
  exists it is less than~$s$ itself.
\end{proof}

\begin{prop}
  The function $\fn{append}(s,a)$, which returns the result of appending $a$ to
  the sequence $s$, is primitive recursive.
\end{prop}

\begin{proof}
  $\fn{append}$ can be defined by:
  \[
  \fn{append}(s,a) =
  \begin{cases}
    2^{a+1} & \text{if $s = 0$ or $s = 1$} \\
    s \cdot p_{\len{s}}^{a+1} & \text{otherwise.}
  \end{cases}
  \]
\end{proof}

\begin{prop}
  The function $\fn{element}(s,i)$, which returns the $i$th element of $s$
  (where the initial element is called the $0$th), or $0$ if $i$ is
  greater than or equal to the length of $s$, is primitive recursive.
\end{prop}

\begin{proof}
Note that $a$ is the $i$th element of~$s$ iff $p_i^{a+1}$ is
the largest power of~$p_i$ that divides~$s$, i.e., $p_i^{a+1} \mid s$
but $p_i^{a+2} \nmid s$. So:
  \[
  \fn{element}(s,i) =
  \begin{cases}
    0 & \mbox{if $i \geq \len{s}$} \\
    \bmin{a < s}{(p_i^{a+2} \nmid s)} & \text{otherwise.}
  \end{cases}
  \]
\end{proof}

Instead of using the official names for the functions defined above,
we introduce a more compact notation. We will use $(s)_i$ instead of
$\fn{element}(s,i)$, and $\tuple{s_0, \dots, s_k}$ to abbreviate
\[
\fn{append}(\fn{append}(\dots \fn{append}(\emptyseq,s_0)
\dots),s_k).
\]
Note that if $s$ has length~$k$, the elements of $s$ are
$(s)_0$, \dots,~$(s)_{k-1}$.

\begin{prop}
The function $\fn{concat}(s,t)$, which concatenates two
sequences, is primitive recursive.
\end{prop}

\begin{proof}
  We want a function $\fn{concat}$ with the property that
  \[
    \fn{concat}(\tuple{a_0, \dots, a_k}, \tuple{b_0, \dots, b_l}) =
    \tuple{a_0, \dots, a_k, b_0, \dots, b_l}.
  \]
  We'll use a ``helper'' function
  $\fn{hconcat}(s,t,n)$ which concatenates the first $n$ symbols of $t$
  to~$s$. This function can be defined by primitive recursion as
  follows:
  \begin{align*}
    \fn{hconcat}(s,t,0) & = s\\
    \fn{hconcat}(s,t,n+1) & = \fn{append}(\fn{hconcat}(s,t,n),(t)_n)
    \intertext{Then we can define $\fn{concat}$ by}
    \fn{concat}(s,t) & = \fn{hconcat}(s,t,\len{t}).
  \end{align*}
\end{proof}

We will write $s \concat t$ instead of $\fn{concat}(s,t)$.

It will be useful for us to be able to bound the numeric code of a
sequence in terms of its length and its largest element. Suppose $s$
is a sequence of length~$k$, each element of which is less than or equal
to some number~$x$. Then $s$ has at most $k$ prime factors, each at
most~$p_{k-1}$, and each raised to at most $x+1$ in the prime
factorization of~$s$. In other words, if we define
\[
\fn{sequenceBound}(x,k) = p_{k-1}^{k \cdot (x+1)},
\]
then the numeric code of the sequence~$s$ described above is at
most~$\fn{sequenceBound}(x,k)$.

Having such a bound on sequences gives us a way of defining new
functions using bounded search. For example, we can define
$\fn{concat}$ using bounded search. All we need to do is write down a
primitive recursive \emph{specification} of the object (number of the
concatenated sequence) we are looking for, and a bound on how far to
look. The following works:
\begin{align*}
  \fn{concat}(s,t) = {} & \bmin{v < \fn{sequenceBound}(s+t,\len{s} +
    \len{t})}{} \\
  & \quad(\len{v} = \len{s} + \len{t} \land {}\\
    & \qquad \bforall{i < \len{s}}{((v)_i = (s)_i) \land {} \\
      & \qquad \bforall{j < \len{t}}{((v)_{\len{s}+j} = (t)_j)})}
\end{align*}

\begin{prob}
Show that there is a primitive recursive function~$\fn{sconcat}(s)$
with the property that
\[
\fn{sconcat}(\tuple{s_0, \dots, s_k}) = s_0 \concat \dots \concat s_k.
\]
\end{prob}

\begin{prob}
Show that there is a primitive recursive function~$\fn{tail}(s)$
with the property that
\begin{align*}
  \fn{tail}(\emptyseq) & = 0 \text{ and}\\
  \fn{tail}(\tuple{s_0, \dots, s_{k}}) & = \tuple{s_1, \dots, s_{k}}.
\end{align*}
\end{prob}

\begin{prop}
  \ollabel{prop:subseq}
  The function $\fn{subseq}(s, i, n)$ which returns the subsequence
  of $s$ of length~$n$ beginning at the $i$th element, is primitive
  recursive.
\end{prop}

\begin{proof}
  Exercise.
\end{proof}

\begin{prob}
  Prove \olref[cmp][rec][seq]{prop:subseq}.
\end{prob}
  


% ==== OLP file: content/computability/recursive-functions/trees.tex ====
% Part: computability
% Chapter: recursive-functions
% Section: trees



\olfileid{cmp}{rec}{tre}
\olsection{Trees}

Sometimes it is useful to represent trees as natural numbers, just
like we can represent sequences by numbers and properties of and
operations on them by primitive recursive relations and functions on
their codes.  We'll use sequences and their codes to do this. A tree
can be either a single node (possibly with a label) or else a node
(possibly with a label) connected to a number of subtrees. The node is
called the \emph{root} of the tree, and the subtrees it is connected
to its \emph{immediate subtrees}.

We code trees recursively as a sequence $\tuple{k, d_1, \dots, d_k}$,
where $k$ is the number of immediate subtrees and $d_1$, \dots,~$d_k$
the codes of the immediate subtrees. If the nodes have labels, they
can be included after the immediate subtrees. So a tree consisting
just of a single node with label~$l$ would be coded by $\tuple{0,l}$,
and a tree consisting of a root (labelled~$l_1$) connected to two
single nodes (labelled $l_2$, $l_3$) would be coded by $\tuple{2,
  \tuple{0, l_2}, \tuple{0, l_3}, l_1}$.

\begin{prop}
  \ollabel{prop:subtreeseq}
  The function $\fn{SubtreeSeq}(t)$, which returns the code of a
  sequence the elements of which are the codes of all subtrees of the
  tree with code~$t$, is primitive recursive.
\end{prop}

\begin{proof}
  First note that $\fn{ISubtrees}(t) = \fn{subseq}(t, 1, (t)_0)$ is
  primitive recursive and returns the codes of the immediate subtrees
  of a tree~$t$. Now we can define a helper function
  $\fn{hSubtreeSeq}(t,n)$ which computes the sequence of all subtrees
  which are $n$ nodes removed from the root. The sequence of subtrees
  of~$t$ which is $0$ nodes removed from the root---in other words,
  begins at the root of~$t$---is the sequence consisting just
  of~$t$. To obtain a sequence of all level~$n+1$ subtrees of $t$, we
  concatenate the level $n$ subtrees with a sequence consisting of all
  immediate subtrees of the level $n$ subtrees. To get a list of all
  these, note that if $f(x)$ is a primitive recursive function
  returning codes of sequences, then $g_f(s, k) = f((s)_0) \concat
  \dots \concat f((s)_k)$ is also primitive recursive:
    \begin{align*}
      g(s, 0) & = f((s)_0)\\
      g(s, k+1) & = g(s, k) \concat f((s)_{k+1})
    \end{align*}
    For instance, if $s$ is a sequence of trees, then
    $h(s) = g_{\fn{ISubtrees}}(s, \len{s})$ gives the sequence of the
    immediate subtrees of the elements of~$s$. We can use it to define
    $\fn{hSubtreeSeq}$ by
    \begin{align*}
      \fn{hSubtreeSeq}(t, 0) & = \tuple{t} \\
      \fn{hSubtreeSeq}(t, n+1) & = \fn{hSubtreeSeq}(t, n) \concat
      h(\fn{hSubtreeSeq}(t, n)).
    \end{align*}
    The maximum level of subtrees in a tree coded by~$t$, i.e., the
    maximum distance between the root and a leaf node, is bounded by
    the code~$t$. So a sequence of codes of all subtrees of the tree
    coded by~$t$ is given by $\fn{hSubtreeSeq}(t, t)$.
\end{proof}

\begin{prob}
  The definition of $\fn{hSubtreeSeq}$ in the proof of
  \olref[cmp][rec][tre]{prop:subtreeseq} in general includes
  repetitions. Give an alternative definition which guarantees that
  the code of a subtree occurs only once in the resulting list.
\end{prob}



\end{document}
